\section{Vetores}

Um vetor é um conjunto de objetos de mesmo tipo\footnote{Isso porque dois dados do
mesmo tipo possuem o mesmo número de bytes na memória.} que são armazenados contiguamente na
memória.
Seus elementos podem ser acessados diretamente pelos seus índices.
Se o nome do vetor é $S$, então acessamos seu $i$-ésimo membro como $S[i]$.
O nome do vetor é um ponteiro para o seu primeiro elemento, e, ao somá-lo com um
valor inteiro não negativo $i$, obtemos um ponteiro para o elemento de índice $i$ do
vetor, portanto,
\[
    S[i] = *(S + i)
\]

Na sua formulação mais simples, vetores possuem dois atributos principais: tamanho --
o número máximo de elementos que ele suporta -- e o número efetivo de membros.
Chamamos de estáticos os vetores que possuem tamanho previamente determinado e fixo,
enquanto são denominados dinâmicos aqueles cujo tamanho varia conforme a necessidade
da aplicação.
Nas implementações expostas, presume-se que os vetores são dinâmicos.

\subsection{Vetores Não Ordenados}

A busca em vetores não ordenados é realizada percorrendo o vetor sequencialmente até
achar o elemento com a chave desejada.
No pior caso, o elemento ou está na última posição, ou não existe no vetor.
Com isso, a sua complexidade é $O(n)$.
Para vetores, é conveniente retornar o índice do elemento sendo buscado, e, caso ele
não exista, retornamos $-1$ no lugar.

\begin{function}[H]
\SetAlgoLined
\caption{Busca-Vetor($S$, $k$)}
$i \gets 0$\;
\Enqto{$i < S.\NumElem\ \E\ S[i].\Chave \neq k$}{
    $i \gets i + 1$\;
}
\Se{$i = S.\NumElem$}{
    \Retorna{$-1$}\;
}
\Retorna{$i$}\;
\end{function}

A inserção num vetor pode receber como argumentos um ponteiro para o objeto a ser inserido
e o índice da posição em que ele será colocado.
Como pode ser necessário mover os membros já presentes para inserir um novo, ou realocar
o vetor quando esse já está cheio, a inserção tem complexidade $O(n)$.

\begin{function}[H]
\SetAlgoLined
\caption{Insere-Vetor($S$, $x$, $i$)}
\Se{$i < 0\ \Ou\ i > S.\NumElem$}{
   \textbf{erro}\;
}
\Se{$S.\NumElem = S.\Len$}{
    $S \gets \RealocaVetor{$S$}$\;
}
$j \gets S.\NumElem$\;
\Enqto{$j \neq i$}{
    $S[j] \gets S[j-1]$\;
    $j \gets j - 1$\;
}
$S[i] \gets *x$\;
$S.\NumElem \gets S.\NumElem + 1$\;
\end{function}

Contudo, se o vetor for um conjunto não ordenado, então podemos sempre inserir o novo
elemento logo após o último membro, de forma a não executar o laço que move elementos a
partir do índice $i$ uma posição adiante.
Com isso, a inserção assume complexidade $O(1)$ a menos de uma realocação com custo $O(n)$.

\begin{function}
\SetAlgoLined
\caption{Realoca-Vetor($S$)}
    $T \gets \Aloca{$S.\Len * 2$}$\;
    \Para{$i \gets 0\ \Ate\ S.\NumElem - 1$}{
        $T[i] \gets S[i]$\;
    }
    $T.\NumElem \gets S.\NumElem$\;
    $T.\Len \gets S.\Len \cdot 2$\;
    \Retorna{$T$}\;
\end{function}

% Explicação sobre amortização da inserção

Em vetores não ordenados, a remoção pode ser realizada simplesmente copiando o último
elemento para a posição daquele que está sendo retirado.
Desse modo, a remoção é executada em tempo constante.

\begin{function}[H]
\SetAlgoLined
\caption{Remove-Vetor($S$, $i$)}
\Se{$S.\NumElem = 0$}{
    \textbf{erro} ``underflow''\;
}
\Se{$i < 0\ \Ou\ i \ge S.\NumElem$}{
   \textbf{erro}\;
}
$S[i] \gets S[S.\NumElem - 1]$\;
$S.\NumElem \gets S.\NumElem - 1$\;
\end{function}

\subsection{Vetores Ordenados}

Em vetores ordenados, geralmente em ordem crescente, a busca pode ser feita em tempo
sublinear $O(\log n)$ se usarmos o algoritmo de busca binária, que será apresentado num
capítulo mais adiante.
Com ela, em vez de retornar $-1$ quando a chave não é encontrada, retornamos o índice que
o elemento com a chave buscada ocuparia caso estivesse no vetor.
Na inserção, realiza-se a busca pela chave do elemento a ser inserido no vetor para obter
o índice no qual ele deve ser posicionado.
Depois, passa-se esse índice para o mesmo procedimento de inserção apresentado para
vetores não ordenados.

Remoções em vetores ordenados requerem que todo membro que sucedia aquele que está sendo
retirado seja movido para uma posição atrás.
Por consequência, a complexidade da operação é $O(n)$.

\begin{function}[H]
\SetAlgoLined
\caption{Remove-Vetor-Ordenado($S$, $i$)}
\Se{$S.\NumElem = 0$}{
    \textbf{erro} ``underflow''\;
}
\Se{$i < 0\ \Ou\ i \ge S.\NumElem$}{
   \textbf{erro}\;
}
$j \gets i$\;
\Enqto{$j < S.\NumElem - 1$}{
    $S[j] \gets S[j+1]$\;
    $j \gets j+1$\;
}
$S.\NumElem \gets S.\NumElem - 1$\;
\end{function}

Para as operações de consulta de máximo e mínimo, basta retornar o ponteiro para o último
e o primeiro elemento no vetor respectivamente, sendo, pois, triviais por meio da indexação.
O mesmo vale para o predecessor e sucessor do elemento de índice $i$, que consistem nos
elementos de índices $i-1$ e $i+1$ respectivamente quando esses são válidos, senão,
$\nil$ deve ser retornado.
